\documentclass[11pt]{article}
\usepackage[T1]{fontenc}
\usepackage[utf8]{inputenc}
\usepackage{lmodern,amsmath,amssymb,amsthm,mathtools}
\usepackage[margin=0.92in]{geometry}
\usepackage{booktabs,longtable,array,microtype}
\usepackage[dvipsnames]{xcolor}
\usepackage{enumitem,fancyhdr,xurl,hyperref}
\hypersetup{colorlinks=true,linkcolor=MidnightBlue,urlcolor=MidnightBlue,citecolor=MidnightBlue,
 pdftitle={Integral Projections and Lattice Forms for Diophantine Consistency},
 pdfauthor={Prepared for Vladimir Reshetnikov with AI assistance}}
\newtheorem{theorem}{Theorem}[section]
\newtheorem{lemma}[theorem]{Lemma}
\newtheorem{proposition}[theorem]{Proposition}
\newtheorem{corollary}[theorem]{Corollary}
\theoremstyle{remark}\newtheorem{remark}[theorem]{Remark}
\newcommand{\Z}{\mathbb Z}
\newcommand{\Q}{\mathbb Q}
\newcommand{\R}{\mathbb R}
\newcommand{\N}{\mathbb N}
\newcommand{\Con}{\operatorname{Con}}
\newcommand{\htp}{\operatorname{ht}}
\newcommand{\supp}{\operatorname{supp}}
\newcommand{\ZFC}{\mathrm{ZFC}}
\newcommand{\code}[1]{\texttt{\detokenize{#1}}}
\setlength{\emergencystretch}{2em}
\setlength{\headheight}{14pt}
\pagestyle{fancy}\fancyhf{}
\fancyhead[L]{\small Diophantine consistency}
\fancyhead[R]{\small Integral projection and lattice forms}
\fancyfoot[C]{\thepage}
\allowdisplaybreaks[3]
\setlist{itemsep=3pt,topsep=5pt}

\begin{document}
\begin{titlepage}
\vspace*{0.8in}
{\color{MidnightBlue}\Large ProveIt research contribution\par}
\vspace{0.35in}
{\huge\bfseries Integral Projections and Lattice Forms\\[0.15in]
for Diophantine Consistency\par}
\vspace{0.3in}
{\Large Complete explicit quartics and exact certificates\par}
\vspace{0.6in}
{\large Prepared for Vladimir Reshetnikov\\[0.08in]
25 October 2026\\[0.08in]With AI assistance\par}
\vspace{0.5in}
\begin{abstract}
We construct and supply several ordinary integer-coefficient polynomials
whose integer unsolvability is provably equivalent in ZFC to its syntactic
consistency. The contribution is an exact compression of the pinned
ProveIt direct proof-calculus constructions, using two complementary
operations. First, an affine equation with a nonunit coefficient can
eliminate a variable when a retained equation certifies that the rational
restoration is integral. Second, a positive-definite integral Gram matrix
with even diagonal can replace the diagonal sum-of-squares combiner.
The principal result is a completely specified quartic in 240 unrestricted
integer variables with 5,803 collected monomials and coefficient height
6,404. Its Gram form has only thirteen mixed residual products; a rank-eight
block is an even unimodular lattice. A companion has 9,357 monomials and
height 4,160. A simpler sum-of-squares companion has 9,217 monomials and
height 4,512, improving both support and height over the inspected
240-variable predecessor's 21,482 and 6,144.

The paper prints every residual of the principal polynomial. The package
also contains complete collected coefficient arrays, rational restoration
certificates, exact positive-definiteness certificates, independent
checkers, and immutable copies of the inherited source archives.
The new transformations are proved here; the C222-to-ZFC calibration is
explicitly inherited. We do not claim a global simplicity optimum,
literature-wide priority, or a new end-to-end proof-assistant derivation.
An additional, fully materialized eleven-axiom finite theory gives a
precise research direction without being counted as a new compiled
Diophantine endpoint.
\end{abstract}
\vfill
\noindent\small Comparison snapshot:
\code{7a8da414376a80b90ecf799f37ceb26fb9ac674e}.\\
All unknowns in the completed new equations range over $\Z$.
\end{titlepage}

\tableofcontents
\newpage

\section{What is obtained}

The requested logical target is
\begin{equation}\label{eq:target}
\ZFC\vdash
\left[\neg\exists\boldsymbol x\in\Z^n\;P(\boldsymbol x)=0\right]
\longleftrightarrow \Con(\ZFC).
\end{equation}
Here consistency means the absence of a finite refutation in the
repository's specified syntactic ZFC calculus. The equivalence does not
assume consistency and does not prove either side separately.

The central completed polynomial in this report is
\begin{equation}\label{eq:principal}
\boxed{\quad Q_{240}(\boldsymbol x)=
\sum_{i=0}^{176}r_i(\boldsymbol x)^2+
\sum_{(i,j,\epsilon)\in E}\epsilon\,r_i(\boldsymbol x)r_j(\boldsymbol x)=0,
\qquad \boldsymbol x\in\Z^{240}.\quad}
\end{equation}
Every $r_i$ is a fixed quadratic integer polynomial, printed in
Appendix~\ref{app:residuals}. The thirteen-element signed edge list $E$ is
printed in Section~\ref{sec:gram}. Thus the displayed specification is a
finite ordinary polynomial, with no variable exponent, factorial, binomial
primitive, uninstantiated universal polynomial, or undecidable coefficient
selection.

\subsection{Exact measurements}

The support count is the number of nonzero monomials after full collection;
height is the maximum absolute collected coefficient. These are different
from residual count, expression length, operation count, and coefficient
bit complexity. All rows below have exact degree four and all their
declared coordinates occur in the collected polynomial.

%%RESULTS_TABLE%%

The usual sum of 177 residual squares gives two useful projected forms.
The Gram versions retain those residuals and add respectively thirteen and
seven mixed products. They have the same zero sets as their respective
sum-of-squares inputs.

The new sum-of-squares form $S_{240}^{\rm low}$ improves the inspected
240-variable Shared Factorial Error endpoint in both support and height:
\[
21\,482\longrightarrow9\,217,\qquad
6\,144\longrightarrow4\,512.
\]
The changes are reductions of approximately $57.09\%$ and $26.56\%$.
The lower-height Gram companion $H_{240}$ reaches $4\,160$, while its
support rises from $9\,217$ to $9\,357$. The principal $Q_{240}$ uses fewer
terms, $5\,803$, at height $6\,404$. These are distinct tradeoffs; no single
one dominates every other profile.

\subsection{Scope of the comparison}

The pinned repository contains constructions with far fewer unknowns,
including nine-positive-variable forms and a 37-integer quartic, using
different coefficient architectures and circuit costs. The present work
improves the direct numerical proof-calculus family with fully materialized
small integer coefficients. It is not a variable-count record across all
representations. The meaningful novelty here is the certified use of
integrality-aware projection together with a nondiagonal residual
combiner on these complete equations.

The immediate comparison archives are
\code{ConZFC_Shared_Factorial_Error_2026-10-25.zip} and
\code{ConZFC_Factorial_Free_Binomial_CRT_2026-10-25.zip}, both in
\code{docs/incoming} at the stated commit~\cite{repo,shared,factorialfree}.
We inspected all seven incoming archives and the newest consolidated
construction directories. Their earlier source histories are retained in
the original archives rather than silently replaced.

\section{The exact logical dependency}\label{sec:dependency}

Let $P_{\rm src}$ denote one of the supplied predecessor polynomials.
The source reports establish, by their ordinary arithmetic proofs and
identified calibration interfaces,
\begin{equation}\label{eq:source-theorem}
\ZFC\vdash
\left[\exists\boldsymbol z\in\Z^m\;P_{\rm src}(\boldsymbol z)=0\right]
\longleftrightarrow
\left[\text{there is a finite ZFC refutation}\right].
\end{equation}
Their direct arithmetic route encodes derivations in the numerical
calculus C222; it does not simulate K215 or instantiate Jones's universal
equation. Its named logical interfaces are
\begin{align*}
&\code{KEndpoint.c222_iff_zf_refutation},\\
&\code{ConstructibleChoice.zf_refutation_iff_zfc_refutation}.
\end{align*}
The relevant source files are included for inspection and pinned to the
same repository snapshot~\cite{endpoint,choice}. The second interface is
the constructible-universe equiconsistency transport from ZF to ZFC.

The new work proves explicit equivalences
\[
\exists P_{\rm new}=0\ \longleftrightarrow\ \exists P_{\rm src}=0
\]
and, for a Gram companion, equality of the corresponding integer zero
sets. Composing with \eqref{eq:source-theorem} and negating gives
\eqref{eq:target}. Every new argument uses elementary arithmetic of
integers and rationals, finite polynomial identities, or an exact finite
positive-definiteness certificate; all are formalizable in ZFC.

This is an ordinary mathematical proof composed with a named inherited
theorem. The release does not contain a newly checked object-language
ZFC derivation of the entire chain, nor a rebuilt transitive Lean/Rocq
dependency closure. Finite coefficient checks authenticate the
polynomials and transformations; they do not substitute for the
unbounded source interpretation. The full predecessor reports and source
proof closures are retained in the two original archives under
\code{sources/}.

\section{Projection through a nonunit coefficient}\label{sec:projection}

\subsection{The general theorem}

Ordinary gate elimination usually restricts a defining equation to
$x=f(\boldsymbol y)$ or $-x=f(\boldsymbol y)$. A coefficient such as $2$
is excluded because solving $2x=f(\boldsymbol y)$ over $\Q$ can create
spurious integer tuples $\boldsymbol y$ with a half-integral $x$.
The missing condition can sometimes be supplied by a different equation
that is already present.

\begin{theorem}[Integral restoration]\label{thm:projection}
Let a finite integer-polynomial system in $(x,\boldsymbol y)$ contain
\[
a x+L(\boldsymbol y)=0,\qquad 0\ne a\in\Z,
\]
and a different retained equation
\[
b(\boldsymbol y)x^d+
h_{d-1}(\boldsymbol y)x^{d-1}+\cdots+h_0(\boldsymbol y)=0,
\qquad d\ge1,
\]
with integer-polynomial coefficients. Assume
$\gcd(a,b(\boldsymbol y))=1$ for each integer tuple under consideration.
Substitute $x=-L(\boldsymbol y)/a$ into every other equation, clear fixed
rational denominators by nonzero constants, and delete the first equation.
Then coordinate projection is a bijection between the old and new
integer zero sets. The inverse inserts $x=-L(\boldsymbol y)/a$.
If $L$ is affine and every old residual has degree at most two, every new
residual still has degree at most two.
\end{theorem}

\begin{proof}
The forward map follows by substitution. For the reverse map, start with
an integer zero $\boldsymbol y$ of the new system and set
$x=-L(\boldsymbol y)/a\in\Q$. Every old equation vanishes, because its
cleared version vanishes and each clearing constant is nonzero.
Write $x=u/v$ in lowest terms, with $v>0$. The affine equation implies
$v\mid a$. Multiplying the retained degree-$d$ equation by $v^d$ shows
$v\mid b(\boldsymbol y)u^d$. Since $\gcd(u,v)=1$, one has
$v\mid b(\boldsymbol y)$. The stated coprimality forces $v=1$.
Thus $x$ is integral. The inverse is unique because $a\ne0$.
Finally, affine substitution does not increase the total degree of a
polynomial. Removing the nonzero integer content of a residual also
preserves its zero set.
\end{proof}

A monic retained equation, or one with leading coefficient $-1$, is
enough. The proof is the standard rational-root argument. The contribution
is the verified choice of pivots and integrality equations in the
complete source systems. Without the second equation, deleting $x$ from
$2x-y=0$ would incorrectly permit odd $y$.

\subsection{The selected square witness}

Both selected Shared Factorial Error parents have 241 unrestricted
coordinates $z_0,\ldots,z_{240}$ and 178 residuals. Use zero-based row and
variable indices. Source row 88 is
\[
A(\boldsymbol z)-2z_{39}=0,
\]
where the affine expression independent of $z_{39}$ is
\begin{align}
A(\boldsymbol z)={}&z_{103}-z_{10}-z_{11}
 -4(z_{21}+z_{26}+z_{32}+z_{38}+z_{67})\notag\\
&-2(z_{24}+z_{29}+z_{31}+z_{34}+z_{36}
      +z_{42}+z_{46}+z_{68}).\label{eq:A}
\end{align}
The different source row 23 is
\begin{equation}\label{eq:square-row}
z_{39}-(z_{11}+2z_{38})^2=0.
\end{equation}
Delete row 88 and substitute $z_{39}=A/2$ in every other row.
After denominator clearing, \eqref{eq:square-row} becomes
\begin{equation}\label{eq:integrality-row}
A-2(z_{11}+2z_{38})^2=0.
\end{equation}
At any integer zero, $A$ is therefore even and the rationally reconstructed
$A/2$ is exactly the integer square $(z_{11}+2z_{38})^2$.
No new parity coordinate is needed.

All other residuals stay quadratic. Reindex retained coordinates in their
old order:
\[
x_j=z_j\quad(j<39),\qquad x_j=z_{j+1}\quad(j\ge39).
\]
The certificate records each nonzero rational multiplier used to clear
denominators and remove residual content. The projected balanced source
has 5,477 collected monomials and height 15,107; the projected lower-height
source has 9,217 and 4,512.

Eliminating $z_{103}$ through the unit coefficient in the same row is a
different projection. It expands many occurrences of a large sum.
Eliminating the square coordinate $z_{39}$ instead uses more arithmetic
information but produces a much sparser 240-variable endpoint.
This explains why a rule that considers only unit pivots can miss a
useful dimension-preserving comparison.

\subsection{A second application}

The factorial-free principal source has 242 coordinates and 178 rows.
Its row 88 has coefficient $-2$ at $z_{67}$, and retained row 52 says
\[
z_{67}=2+4z_{66}+4z_{66}^2.
\]
The same argument removes $z_{67}$ and gives an explicit 241-variable
quartic with 177 residuals, 5,141 monomials, and height 7,968.
The predecessor's other 241-variable projection has 12,774 terms at
the same height. The complete new system and every substitution are
included in \code{projection/factorialfree241/}.

\section{A positive lattice form in place of a sum of squares}\label{sec:gram}

\subsection{The aggregation theorem}

\begin{theorem}[Even integral Gram aggregation]\label{thm:gram}
Let $r_0,\ldots,r_{m-1}\in\Z[\boldsymbol x]$ and let $H$ be a symmetric,
positive-definite integer matrix with even diagonal. Then
\[
Q(\boldsymbol x)=\frac12\,\boldsymbol r(\boldsymbol x)^T
H\boldsymbol r(\boldsymbol x)
\]
has integer coefficients and, for every real tuple $\boldsymbol x$,
\[
Q(\boldsymbol x)=0
\quad\Longleftrightarrow\quad
r_i(\boldsymbol x)=0\ \hbox{for every }i.
\]
If each residual has degree at most two, then $\deg Q\le4$.
\end{theorem}

\begin{proof}
The diagonal coefficient is $H_{ii}/2\in\Z$; each mixed coefficient is
$H_{ij}\in\Z$, because the two symmetric entries cancel the factor $1/2$.
Positive definiteness gives $\boldsymbol v^TH\boldsymbol v>0$ for every
nonzero real vector $\boldsymbol v$. Apply this to
$\boldsymbol v=\boldsymbol r(\boldsymbol x)$. The degree bound follows
by multiplying two residuals.
\end{proof}

The familiar sum of squares uses $H=2I$. There is no need for that
particular matrix if the only purpose of the final polynomial is to
enforce simultaneous vanishing. Small off-diagonal coefficients can cancel
large coefficients in the expanded polynomial.

This enlarges the aggregation family beyond integer changes of the
residual row basis followed by an unweighted sum of squares. In formal
residual coordinates, an integer sum of squares of integer linear forms
has even mixed coefficients. Our selected form has odd mixed coefficients.
This observation concerns the formal residual variables; it does not
assert that the substituted polynomial admits no other sum-of-squares
representation.

\subsection{The complete thirteen-edge specification}

Number the projected balanced residuals from $0$ to $176$. In addition to
their ordinary squares, the principal polynomial contains
\begin{align}
&-r_{12}r_{16}-r_{12}r_{19}-r_{12}r_{17}
 -r_{49}r_{51}-r_{19}r_{51}\notag\\
&-r_6r_{21}-r_{17}r_{26}-r_{15}r_{21}
 +r_{14}r_{21}+r_{17}r_{19}\notag\\
&-r_7r_{19}-r_7r_{17}+r_6r_{50}.\label{eq:edges}
\end{align}
Equations \eqref{eq:principal}, \eqref{eq:edges}, and
Appendix~\ref{app:residuals} specify every coefficient of $Q_{240}$.

Only two blocks of $H$ are nontrivial. On residual coordinates
$(6,14,15,21,50)$ the first is
\[
H_5=
\begin{pmatrix}
2&0&0&-1&1\\
0&2&0&1&0\\
0&0&2&-1&0\\
-1&1&-1&2&0\\
1&0&0&0&2
\end{pmatrix}.
\]
On coordinates $(7,12,16,17,19,26,49,51)$ the second is
\[
H_8=
\begin{pmatrix}
2&0&0&-1&-1&0&0&0\\
0&2&-1&-1&-1&0&0&0\\
0&-1&2&0&0&0&0&0\\
-1&-1&0&2&1&-1&0&0\\
-1&-1&0&1&2&0&0&-1\\
0&0&0&-1&0&2&0&0\\
0&0&0&0&0&0&2&-1\\
0&0&0&0&-1&0&-1&2
\end{pmatrix}.
\]
The other 164 blocks are $(2)$. Exact rational $LDL^T$ elimination, in the
displayed orders, gives
\[
D_5=\operatorname{diag}(2,2,2,\tfrac12,1),\qquad
D_8=\operatorname{diag}(2,2,\tfrac32,\tfrac56,
                         \tfrac45,\tfrac34,2,\tfrac16).
\]
Every pivot is positive. The full unit lower triangular factors are
supplied as rational numerator--denominator pairs in
\code{gram/balanced/GRAM_CERTIFICATE.json}; the checker verifies
$H=LDL^T$ entry by entry. This is an exact proof, independent of any
floating-point eigenvalue estimate.

The two determinants are $4$ and $1$. The second block is an even,
positive-definite, unimodular lattice of rank eight, so it is isometric
to $E_8$ by the classical uniqueness theorem; Elkies gives a short
author-hosted proof~\cite{elkies}. This identification is interesting but
is not needed for the zero-set equivalence. The rational $LDL^T$
certificate already proves everything required.

\subsection{Measured effect and companion}

The Gram change takes the balanced projected input from height 15,107 to
6,404, a reduction of $57.61\%$. Its support grows from 5,477 to 5,803,
or $5.95\%$. The support/height tradeoff is explicit and does not change
the dimension or degree.

The lower-height input has a separate seven-edge Gram form, supplied
completely in \code{gram/lowheight/}. It changes 9,217 terms and height
4,512 to 9,357 and 4,160. Its two nontrivial blocks have orders seven
and two, determinants $4$ and $3$, and exact positive pivots
\[
(2,2,\tfrac32,2,1,\tfrac43,\tfrac14),
\qquad (2,\tfrac32).
\]
That form is $H_{240}$ in the results table.

\begin{corollary}[Completed consistency equations]
The principal polynomial $Q_{240}$, the lower-height companion $H_{240}$,
both projected sum-of-squares polynomials, and the factorial-free
241-variable companion each satisfy \eqref{eq:target} in their stated
unrestricted integer domains.
\end{corollary}

\begin{proof}
Theorem~\ref{thm:projection} transports integer solutions to the exact
source equation. For either Gram companion,
Theorem~\ref{thm:gram} first transports zeros to all its residuals.
Apply the inherited source equivalence \eqref{eq:source-theorem} and
negate the resulting existential statement.
\end{proof}

\section{Why the new equation is fully explicit}

The main JSON record contains exactly 240 ordered variables and 177
quadratic residual lists. A sparse term
\[
[c,[i,j,\ldots]]
\]
denotes $c x_i x_j\cdots$; repeated indices denote repeated factors and
an empty list denotes the constant $c$. All indices are zero-based.
The additional list \code{gram_cross_terms} records the triples
$(i,j,\epsilon)$ in \eqref{eq:edges}.

\code{expanded_polynomial.json} contains every collected coefficient of
the single ordinary quartic. The specification does not depend on a
search run, an optimizer, a positivity oracle, or an unexpanded
universal-polynomial parameter. Evaluators use integer arithmetic only.
The fact that the equation encodes an undecidable consistency question
does not prevent its coefficients, dimensions, degree, and evaluation
at any given finite integer tuple from being calculated exactly.

The transformations also retain more than bare satisfiability.
Projection and affine restoration are mutually inverse on the old
and new integer zero sets. A Gram change preserves the entire real
zero set of the residual conjunction. These stronger statements make
the soundness argument particularly short and auditable.

\section{Verification and its limits}\label{sec:verification}

The top-level \code{verify.py} checks the exact source-archive bindings,
all new projection identities, both Gram certificates, every final
collected coefficient, and the generated report data. It uses Python's
standard library. The projection checker uses exponent-vector
polynomials and does not import the producer, whose monomials are lists
of indices. The Gram checker independently multiplies residuals,
reconstructs each block, and checks the rational factorization.

The retained receipts include:
\begin{itemize}
\item full coefficient and integer-restoration checks for all three
      projections, including nine deliberately corrupted controls;
\item independent checks of the thirteen- and seven-edge Gram forms,
      including rejection of a singular Gram matrix and a false
      factorization;
\item exact regeneration of the projection JSON data and both Gram
      endpoints;
\item a fresh replay of the inherited Shared Factorial Error release,
      including its parent arithmetic artifact checks and its printed
      equation check;
\item separate reviews of the projection, Gram forms, and the
      inherited signed-radix/nearest-factorial proof boundary;
\item exact evaluation-circuit and formula-tree checks for the
      eleven-axiom finite theory, including all 5,803 collected
      coefficients and five deliberately corrupted controls.
\end{itemize}

Required checks use explicit failures rather than Python assertions, and
normal and optimized-interpreter runs are recorded. Independent reviews
here mean separately performed reviews within an AI-assisted workflow,
not outside human peer review.

\subsection{A useful stopping certificate}

For each projected system, form the matrix whose rows are the
coefficients of the homogeneous quadratic parts of its 177 residuals.
The supplied finite-field elimination gives rank 177 modulo the prime
$1\,000\,003$. A nonzero minor modulo a prime is a nonzero integer
minor, so the rows are independent over $\Q$. Therefore no nonzero
constant rational linear combination of these residuals is affine.

This rules out exposing another affine row by ordinary constant row
combinations in these fixed systems. It does not rule out a different
projection order, polynomial multipliers, a changed source certificate,
or a different proof calculus. It is a narrow and useful stopping result,
not a lower bound on Diophantine representations of consistency.

\subsection{What was not checked}

We did not search for a zero, decide consistency, rebuild all historical
Lean/Rocq dependencies, or prove a global optimum. No finite collection
of sampled tuples would establish the universal logical equivalence.
The checks establish literal algebra and finite certificates. The
ordinary proofs establish the new universal transformations, and the
credited source proofs provide the refutation interpretation.

\section{A tested semantic redesign that did not improve the endpoint}

The source theory can be altered safely by separating its full
membership-vacuity family from a fixed-role seed family. One may also
replace the four-parameter implicational axiom by three-parameter
classical bases and express freshness using finitely many polynomial
charts. However, reducing local schema arity does not by itself
reduce the final Diophantine equation.

We compiled two explicit probes through the portable natural-number
compiler. The three-parameter variant produced 246 natural coordinates
and 201 residuals; the split four-parameter variant produced 226 and 178.
Both exceed the inspected 219-natural-variable endpoint. The extra
branches and constructor work outweigh the anticipated saving.
The repository's incoming Digit Recovery experiments already contain
a related three-parameter probe, so that direction is credited rather
than claimed as unexplored.

These probes are not release endpoints and are not used in the proof of
the completed quartics. Their measured failure is useful evidence for
optimizing the entire compiled certificate, rather than optimizing only
schema counts or the number of local parameters.

\section{A different research route: a finite theory with one coded input}
\label{sec:grounding}

The preceding results are complete. The following construction is a
proved logical transformation and a materialized finite-theory target
for a new compiler; it does not yet supply a smaller materialized
Diophantine equation.

Start with an explicitly given integer polynomial $P$ and a finite
arithmetic circuit that evaluates it. Consider the fixed finite language
of commutative rings with identity, augmented by a binary pairing
function $\pi$ and unary projections $\ell,r$. The eight function symbols
are $0,1,-,\ell,r,+,\cdot,\pi$, with respective arities
$0,0,1,1,1,2,2,2$. Let $E$ contain the universal closures of the following
eight ring identities:
\begin{align*}
(a+b)+c&=a+(b+c), & a+b&=b+a,\\
a+0&=a, & a+(-a)&=0,\\
(ab)c&=a(bc), & ab&=ba,\\
a1&=a, & a(b+c)&=ab+ac,
\end{align*}
and the universal closures of
\[
\ell(\pi(a,b))=a,\qquad r(\pi(a,b))=b.
\]
Use a balanced binary pairing tree to encode a $k$-tuple as one element.
The corresponding coordinate projections have depth at most
$\lceil\log_2 k\rceil$ after padding to a power of two.

If the circuit for $P$ has $s$ stored wires, a second balanced tuple
codes their values. Write $C_P(z,u)$ for the finite conjunction of its
gate equations, where input values are projections of $z$ and wire
values are projections of $u$. Write $o(u)$ for the output projection.
Define the single first-order sentence
\begin{equation}\label{eq:theta}
\theta_P:\quad \forall z\,\exists u\;
\bigl(C_P(z,u)\ \wedge\ o(u)\ne0\bigr).
\end{equation}
All conjunctions can be balanced; the logical form has only the two
displayed outer quantifiers, though the finite matrix and its projection
terms must still be counted.

\begin{theorem}[Finite pairing-ring reduction]\label{thm:ring}
For every explicit integer polynomial $P$ with its fixed finite
arithmetic circuit, ZFC proves
\[
\Con(E+\theta_P)
\quad\Longleftrightarrow\quad
\neg\exists\boldsymbol n\in\Z^k\;P(\boldsymbol n)=0.
\]
\end{theorem}

\begin{proof}
If $P$ has no integer zero, use the standard ring $\Z$. Choose a
bijection $\pi:\Z^2\to\Z$ with its two inverse coordinate functions.
For any $z$, its projections are integers. The circuit has its uniquely
computed integer wire values; pair those values into $u$. Its output
is nonzero, so \eqref{eq:theta} holds. This is a model of the finite
theory, hence the theory is consistent by first-order soundness.

Conversely, suppose an integer tuple $\boldsymbol n$ is a root.
Form its ground paired term $z_{\boldsymbol n}$ using ordinary ring
numerals, including additive inverses for negative integers.
The projection identities reduce its input coordinates to those
numerals. Instantiate \eqref{eq:theta} at $z_{\boldsymbol n}$.
Its gate equations force, by a finite induction on the circuit,
each projected wire value to equal its calculated integer numeral.
Ring numeral normalization proves that the final value is $0$.
This contradicts the required nonzero output. These steps give a
finite first-order refutation, uniformly constructible from the root
and the finite circuit. The two arguments are formalizable in ZFC.
\end{proof}

Applying this theorem to a completed polynomial in this report gives
another explicit finite theory equiconsistent with ZFC. It does not
need a new arithmetized ZFC proof predicate, Robinson arithmetic, or a
finite class-theory presentation.

\subsection{A fully materialized instance}

The directory \code{finite_theory/} supplies this theory for the principal
$Q_{240}$, including its complete arithmetic directed acyclic graph
and all eleven literal first-order syntax trees. Integer constants are
constructed from $0$ and $1$ through binary addition chains and
negation. Every occurrence of a tuple projection is expanded into
compositions of $\ell$ and $r$; every conjunction is a literal balanced
binary tree. There are no unexpanded tuple macros or gate-specific
function symbols in the artifact.

\begin{center}\small
\begin{tabular}{lr}
\toprule
Measurement & Exact value\\\midrule
Arithmetic gates & 2,736\\
Addition / multiplication / negation gates & 1,200 / 1,040 / 254\\
Input gates / zero gates / one gates & 240 / 1 / 1\\
Arithmetic graph depth (edges) & 23\\
Input / gate tuple projection depths (edges) & 8 / 12\\
Syntax nodes in $\theta_{Q_{240}}$ & 107,256\\
Syntax nodes in the eleven-sentence theory & 107,353\\
Maximum functional syntax depth (edges) & 28\\
Formula-only depth / maximum term depth (edges) & 15 / 13\\
Complete theory JSON size (bytes) & 697,036\\
\bottomrule
\end{tabular}
\end{center}

The sentence $\theta_{Q_{240}}$ has exactly two quantifier nodes and
quantifier depth two. Counting the ten background sentences as well,
the theory has twenty-two quantifier nodes and maximum quantifier
depth three. The syntax-node counts include repeated projection terms;
whole-theory counts are sums and depths are maxima over its eleven
sentences. These are measurements of the functional language, not
measurements of a relational translation or a numerical formula code.

An independent checker expands the complete arithmetic graph and,
separately, the residual Gram form, comparing all 5,803 coefficients.
It checks every syntax node, symbol arity, bound variable, background
axiom, gate equation, conjunction, and metric. The builder regenerates
the five mathematical JSON files byte for byte. This supplies a
concrete compilation input rather than merely a schematic sentence.

\subsection{The remaining compilation problem}

The engineering opportunity is to compile refutations of this finite
theory with a fixed logical kernel, moving the particular polynomial
into fixed finite axiom-code data. Balanced tuple coding avoids the naive
prefix of one quantifier per input or intermediate wire. It is not
free compression: a complete implementation must relationalize the
function language, supply the exact axiom codes, account for their
bit lengths, and prove the new source-to-calculus calibration. The
present Cantor-style formula coding can make a deep encoded sentence
astronomically large. No improved variable count, degree, or practical
coefficient size is asserted for this route until those costs are
measured and the final polynomial is emitted.

\section{Files and incorporation}

The package is intended as a new contribution directory; it does not
overwrite an existing numbered construction. The principal files are:
\begin{center}\small
\begin{tabular}{p{0.40\textwidth}p{0.51\textwidth}}
\toprule
Path & Purpose\\\midrule
\code{gram/balanced/} & Principal $Q_{240}$, full coefficients and $LDL^T$ certificate\\
\code{gram/lowheight/} & Lower-height $H_{240}$ with its complete certificate\\
\code{projection/} & Three exact projections, source copies, restoration maps and rank certificates\\
\code{finite_theory/} & Complete eleven-axiom syntax trees, 2,736-gate circuit and independent checker\\
\code{sources/} & Two immutable parent archives and pinned calibration excerpts\\
\code{review/} & Actual verification receipts and scoped independent reviews\\
\code{verify.py} & Offline verification entry point\\
\code{evaluate.py} & Exact evaluation at a supplied integer tuple\\
\code{build_report.py} & Rebuild this standalone LaTeX source and PDF\\
\code{SOURCE_LEDGER.json} & Source archive hashes, commit, paths and provenance\\
\code{MANIFEST.sha256} & Integrity hashes of the delivered payload\\
\bottomrule
\end{tabular}
\end{center}

Use Python 3.10 or later:
\begin{verbatim}
python3 -B verify.py
python3 -B -O verify.py
python3 -B verify.py --rebuild
python3 -B evaluate.py --profile gram-balanced point.json
python3 -B build_report.py
\end{verbatim}
The last command additionally requires a conventional pdfLaTeX
installation with the packages listed in the source preamble.
PDF metadata and TeX-engine versions can change PDF bytes while the
mathematical JSON data remain identical.

\appendix
\section{Every residual of the principal equation}\label{app:residuals}

All $x_0,\ldots,x_{239}$ range independently over $\Z$.
The following finite list defines $r_0,\ldots,r_{176}$.
Together with the thirteen mixed terms in \eqref{eq:edges}, this is
the complete polynomial $Q_{240}$. The formulas are generated from
the canonical JSON and checked against it.

%%RESIDUAL_APPENDIX%%

\section{Machine-readable specification and exact evaluation}

Each equation record declares its domain and ordered coordinates.
For a sum-of-squares profile, evaluate every quadratic residual and add
its square. For a Gram profile, also add each signed mixed product in
\code{gram_cross_terms}. The collected coefficient file is a second
exact representation of the same polynomial. The verifier expands the
residual expression and compares it to that full coefficient list.

No externally supplied semantic witness is needed to evaluate a tuple.
The evaluators are not solvers. A result at any particular tuple gives
no decision procedure for the existence of an arbitrary integer zero.

\begin{thebibliography}{99}
\bibitem{repo}
V. Reshetnikov and contributors, \emph{ProveIt}, 25 October 2026.\\
Snapshot \code{7a8da414376a80b90ecf799f37ceb26fb9ac674e}.
\url{https://github.com/VladimirReshetnikov/ProveIt/tree/7a8da414376a80b90ecf799f37ceb26fb9ac674e/Computability/DiophantineConsistency}.

\bibitem{shared}
ProveIt incoming research package,
\emph{Shared factorial errors: explicit Diophantine equations for Con(ZFC)},
25 October 2026, prepared with AI assistance.
\url{https://github.com/VladimirReshetnikov/ProveIt/blob/7a8da414376a80b90ecf799f37ceb26fb9ac674e/docs/incoming/ConZFC_Shared_Factorial_Error_2026-10-25.zip}.
Original archive retained under \code{sources/}.

\bibitem{factorialfree}
ProveIt incoming research package,
\emph{Removing factorial authenticity from Diophantine consistency},
25 October 2026, prepared with AI assistance.
\url{https://github.com/VladimirReshetnikov/ProveIt/blob/7a8da414376a80b90ecf799f37ceb26fb9ac674e/docs/incoming/ConZFC_Factorial_Free_Binomial_CRT_2026-10-25.zip}.
Original archive retained under \code{sources/}.

\bibitem{endpoint}
ProveIt, \code{KEndpoint.lean}, theorem
\code{c222_iff_zf_refutation}, pinned snapshot above.
\url{https://github.com/VladimirReshetnikov/ProveIt/blob/7a8da414376a80b90ecf799f37ceb26fb9ac674e/Computability/BusyBeaver/Lean/BusyBeaver/ZFC/KEndpoint.lean}.

\bibitem{choice}
ProveIt, \code{ConstructibleChoice.lean}, theorem
\code{zf_refutation_iff_zfc_refutation}, pinned snapshot above.
\url{https://github.com/VladimirReshetnikov/ProveIt/blob/7a8da414376a80b90ecf799f37ceb26fb9ac674e/Computability/BusyBeaver/Lean/BusyBeaver/ZFC/ConstructibleChoice.lean}.

\bibitem{elkies}
N. D. Elkies, \emph{Yet another proof of the uniqueness of the $E_8$
lattice}, August 2004. The stated uniqueness theorem is credited to
L. J. Mordell (1938).
\url{https://people.math.harvard.edu/~elkies/Misc/E8.pdf}.

\bibitem{jones}
J. P. Jones, \emph{Universal Diophantine equation},
Journal of Symbolic Logic \textbf{47} (1982), 549--571.
\url{https://doi.org/10.2307/2273588}.
Background on universal representations, not a dependency of the
direct arithmetic route used here.

\bibitem{bayer}
J. Bayer, M. David, M. Ha{\ss}ler, Y. Matiyasevich and D. Schleicher,
\emph{Diophantine equations over $\Z$: universal bounds and parallel
formalization}, arXiv:2506.20909v3, 22 October 2026.
\url{https://arxiv.org/html/2506.20909v3}.
Background on the distinct variable/degree optimization problem.
\end{thebibliography}
\end{document}
